\section{Conclusão}
\label{sec:conclusao}

Este trabalho apresentou o desenvolvimento do conversor CC-CC de três
portas do tipo VR-BESS para integração de painel fotovoltaico, banco de
baterias e carga, com vistas à futura aplicação da técnica FCS-MPC. Foram
implementados, como blocos mascarados próprios em substituição às
bibliotecas nativas do Simscape Electrical (indisponíveis no ambiente
MATLAB Online utilizado), os modelos do painel fotovoltaico de diodo único
e da bateria de íon-lítio segundo a formulação de Shepherd, integrados ao
domínio físico do conversor por meio de sensores e fontes controladas.

A validação em malha aberta, conduzida com as chaves sob razão cíclica
fixa, confirmou a correta integração entre esses subsistemas e o
conversor: as grandezas de regime permanente observadas em cada porta --
corrente do arranjo fotovoltaico, tensão e corrente de bateria, estado de
carga e tensão de saída -- mostraram-se coerentes com os parâmetros de
projeto adotados (configuração série/paralelo dos arranjos, parâmetros de
catálogo do painel e da bateria), inclusive quanto à variação do estado de
carga frente à corrente de bateria integrada. Esses resultados consolidam
a base de simulação -- planta, sensores e atuadores no domínio físico --
sobre a qual o controlador será desenvolvido.

A validação também evidenciou que a tensão do arranjo fotovoltaico em
regime depende da corrente demandada pelo conversor -- e não apenas da
associação série/paralelo adotada --, de modo que a razão cíclica de
$S_2$ precisou ser calculada a partir do ponto de operação
autoconsistente do arranjo sob a corrente do Modo~1
(Seção~\ref{subsec:duty_cycle}), e não a partir da tensão de máxima
potência nominal. Essa calibração é válida para a corrente de
carregamento de bateria observada nesta validação; variações nessa
corrente -- decorrentes, por exemplo, da evolução do estado de carga da
bateria ou de mudanças na irradiância -- deslocam o ponto de operação e
exigiriam uma nova calibração da razão cíclica, o que reforça a
necessidade do controlador em malha fechada proposto para a etapa
subsequente do trabalho.

O trabalho não abrangeu, até o momento, a implementação da malha de
controle: os resultados aqui apresentados refletem o comportamento do
conversor sob comutação em malha aberta, sem regulação ativa de tensão de
saída ou de corrente de bateria, e para um único ponto de operação, sem
degraus de irradiância, carga ou referência. A etapa subsequente deste
trabalho consiste em implementar o controlador FCS-MPC descrito na
Seção~\ref{subsec:fcsmpc} sobre a planta validada, avaliando sua capacidade
de regulação da tensão de saída e de gerenciamento da corrente de bateria
frente a esses cenários, bem como investigar o emprego dos observadores de
estado apresentados na Seção~\ref{subsec:observadores} como alternativa
para redução do número de sensores físicos necessários à implementação do
controlador.
